\section{Related Work}
\label{sec:related}

\subsection{Water Quality Indices}
Composite indices reduce a vector of measurements to one interpretable number. The
approach originates with Horton~\cite{horton1965} and was formalised in the National
Sanitation Foundation index of Brown et al.~\cite{brown1970wqi}; Abbasi and
Abbasi~\cite{abbasi2012wqi} survey the family that followed. The CCME
index~\cite{ccme2001wqi} combines three factors -- the proportion of variables failing
their guideline (scope), the proportion of individual tests that fail (frequency), and
the magnitude by which they fail (amplitude) -- as a normalised distance from an ideal
state. Lumb et al.~\cite{lumb2011review} review its behaviour and note its sensitivity to
the choice of guideline thresholds and to the temporal window over which frequency is
computed, both of which matter for how the index is interpreted.

\subsection{Machine Learning for Water Quality}
Supervised learning has been applied to water resources along two distinct lines:
predicting unmeasured or delayed constituents from sensor surrogates, and modelling
composite indices directly for automated watershed scoring. Ahmed et
al.~\cite{ahmed2019ml} and Bui et al.~\cite{bui2020hybrid} report strong accuracy for
index prediction with tree ensembles and hybrid models. Our reading of this literature is
that the second line requires care that is not always taken: when the index is a known
function of the model's own inputs, high accuracy is the expected outcome and carries
little information about the model.

We retain the continuous formulation rather than predicting the five CCME classes.
Discretisation discards variance precisely where it matters, since a transition from
79.9 to 80.1 crosses a class boundary while a shift from 50 to 95 within a class is
invisible, and it converts a naturally skewed continuous target into a severely imbalanced
classification problem in which the worst-quality records are under 5\% of the corpus.

\subsection{Tree Ensembles on Tabular Data}
Random Forest~\cite{breiman2001randomforests} reduces variance by averaging decorrelated
trees grown on bootstrap samples with random feature subsets, and rarely overfits as the
ensemble grows. Its impurity-based importances are, however, biased under correlated or
high-cardinality predictors, a caution raised by Strobl et al.~\cite{strobl2007} and
echoed for hydrological applications by Tyralis et al.~\cite{tyralis2019}; we therefore
report SHAP attributions~\cite{lundberg2017shap} rather than impurity importances.
Gradient boosting~\cite{friedman2001gradient} instead fits trees sequentially to the loss
gradient, reducing bias and capturing higher-order interactions, and
XGBoost~\cite{chen2016xgboost} adds regularisation, shrinkage, column subsampling and
sparsity-aware splitting.

Comparative claims between these families are frequently confounded by unequal
experimental conditions -- an extensively tuned boosting model set against a
default-configured forest, or inconsistent partitions between the two. We hold the split,
the preprocessing and the evaluation contract identical across models and availability
regimes, so that the comparison reflects the algorithms rather than the effort spent on
them.

\subsection{Leakage and Optimistic Validation}
Kaufman et al.~\cite{kaufman2012leakage} distinguish leakage in training examples from
leakage in features, where a predictor encodes information about the target that would
not be legitimately available. A target computed by formula from its own predictors is
the extreme case: the mapping is exactly recoverable and generalisation error approaches
the numerical precision of the approximation. Roberts et al.~\cite{roberts2017cv} make
the related argument that structured data demands structured validation, since random
partitioning of spatially or temporally clustered records places near-duplicates on both
sides of the split. Both concerns apply directly here, and shape the formulation in
Section~\ref{sec:method} and the split in Section~\ref{sec:split}.
